% year: תשע"ט
% type: מבחן
% semester: א
% moed: מיוחד
% number: 2
% points: 20
% categories: definite-integrals, antiderivatives
% source: exams/מבחנים/תשעט/מועד מיוחד תשעט.pdf

\begin{question}
חשבו את האינטגרל הבא $\displaystyle\int_0^{\pi/4} x\tan^2 x\,dx$.
\end{question}

\begin{hint}
השתמשו בזהות $\tan^2 x = \frac{1}{\cos^2 x}-1$ ופצלו לשני אינטגרלים.
\end{hint}

\begin{hint}
את $\int x\cdot\frac{1}{\cos^2 x}\,dx$ חשבו באינטגרציה בחלקים, עם $u=x$ ו-$v'=\frac{1}{\cos^2x}$.
\end{hint}

\begin{solution}
לפי הזהות $1+\tan^2 x = \frac{1}{\cos^2 x}$:
\[ \int_0^{\pi/4} x\tan^2 x\,dx = \int_0^{\pi/4} \frac{x}{\cos^2 x}\,dx - \int_0^{\pi/4} x\,dx. \]
(הפונקציות רציפות בקטע $[0,\pi/4]$, ולכן כל האינטגרלים קיימים.)

\textbf{האינטגרל הראשון} — אינטגרציה בחלקים עם $u=x$, $v'=\frac{1}{\cos^2x}$, כלומר $u'=1$, $v=\tan x$:
\[ \int_0^{\pi/4}\frac{x}{\cos^2x}\,dx = \Big[x\tan x\Big]_0^{\pi/4} - \int_0^{\pi/4}\tan x\,dx. \]
מתקיים $\int \tan x\,dx = \int\frac{\sin x}{\cos x}dx = -\ln|\cos x|+C$ (הצבה $t=\cos x$), ולכן
\[ \int_0^{\pi/4}\frac{x}{\cos^2x}\,dx = \frac{\pi}{4}\cdot 1 - 0 + \Big[\ln\cos x\Big]_0^{\pi/4} = \frac{\pi}{4} + \ln\frac{\sqrt2}{2} - \ln 1 = \frac{\pi}{4} - \frac{\ln 2}{2}. \]

\textbf{האינטגרל השני:} $\displaystyle\int_0^{\pi/4}x\,dx = \frac{x^2}{2}\Big|_0^{\pi/4} = \frac{\pi^2}{32}$.

לסיכום:
\[ \int_0^{\pi/4} x\tan^2 x\,dx = \frac{\pi}{4} - \frac{\ln 2}{2} - \frac{\pi^2}{32} \approx 0.1304. \]
\end{solution}
